\documentclass[12pt]{article}
\usepackage[utf8]{inputenc}
\usepackage[T1]{fontenc}
\usepackage{hyperref}
\title{Gamilaraay Language Resources for K-12 Curriculum Development}
\author{Compiled by Claude Sonnet 5.5 (claude-sonnet-5-5) for Andr\'as Kornai}
\date{}
\begin{document}
\maketitle

\section{Introduction}
Gamilaraay (ISO 639-3 kld) is an Australian Pama-Nyungan language of northern New South Wales, in active revival as a second language. It is documented together with Yuwaalaraay and Yuwaalayaay.

\section{Primary Educational Resources}
\begin{itemize}
\item Yuwaalaraay.org teaching and learning page (early learner materials, advanced courses, references): \url{https://yuwaalaraay.com/teaching-and-learning/}
\item Online resources page (picture dictionary \emph{Gaay Garay Dhadhin}, songs CD, children's books): \url{https://yuwaalaraay.com/online-resources/}
\item Garay letters and words literacy resource at Winanga-li: \url{https://www.winanga-li.org.au/language-resource-item/garay-letters-words-literacy}
\item Gamilaraay/Yuwaalaraay page at SALC (Dalang); whether its lessons are still active was not confirmed: \url{https://salc.dalang.com.au/index.php/gamilaraay-yuwaalaraay}
\end{itemize}

\section{Linguistic and Reference Materials}
\begin{itemize}
\item Giacon, development of the Gamilaraay/Yuwaalaraay/Yuwaalayaay dictionary (Sydney University Press): \url{https://open.sydneyuniversitypress.com.au/9781920899554/rl-the-development-of-the-gamilaraay-yuwaalaraay-38-yuwaalayaay-dictionary.html}
\end{itemize}

\section{Religious and Community Resources}
Yuwaalaraay.org and Winanga-li are community-run sites (see above).

\section{Digital and Audio Resources}
\begin{itemize}
\item Endangered Languages Project entry on language learning and teaching: \url{https://www.endangeredlanguages.com/node/109377}
\item Research Data Australia record of recent Gamilaraay/Yuwaalaraay resources: \url{https://researchdata.edu.au/recent-resources-records-gamilaraay-yuwaalaraay/3381231}
\end{itemize}

\section{Curriculum Development Context}
Search results describe a 2003 dictionary (free PDF), an online dictionary, the Gadjigadji app, the multimedia \emph{Gayarragi, Winangali} (over 2,600 audio entries), and university courses. About 60 hours of Yuwaalaraay tapes at AIATSIS underpin the revival. This is a comparatively well-resourced revival language.

\end{document}